\section{Konvergensi Energi dan Optimasi SPSA di Lanskap Kuantum}
\label{sec:4_5_optimasi_spsa}

Proses minimisasi nilai ekspektasi energi $\langle H \rangle$ dilakukan melalui iterasi parameter sirkuit hingga mencapai titik kestabilan global.

\subsection{Mekanisme Estimasi Gradien Stokastik (SPSA)}
\label{subsec:4_5_1_rumus}

Algoritma SPSA digunakan untuk memperbarui parameter $\theta$ dengan mengestimasi gradien fungsi biaya menggunakan perturbasi acak $\Delta$ pada seluruh komponen parameter secara simultan. Secara matematis, estimasi gradien pada iterasi ke-$k$ didefinisikan sebagai:
\begin{equation}
    \hat{g}_k(\hat{\theta}_k) = \frac{y(\hat{\theta}_k + c_k \Delta_k) - y(\hat{\theta}_k - c_k \Delta_k)}{2 c_k} \Delta_k^{-1}
\end{equation}
di mana $y(\cdot)$ adalah hasil evaluasi energi dari sirkuit kuantum dan $c_k$ adalah besarnya perturbasi. Keunggulan utama metode ini adalah efisiensinya yang hanya memerlukan dua evaluasi fungsi per iterasi, sangat cocok untuk sistem kuantum yang memiliki ketidakpastian (\textit{statistical noise}) pada hasil pengukuran.

\subsection{Analisis Sensitivitas \textit{Hyperparameters} terhadap Stabilitas}
\label{subsec:4_5_2_sensitivitas}

Stabilitas konvergensi sangat dipengaruhi oleh pemilihan \textit{hyperparameters} SPSA, terutama faktor penguatan $a$ dan $c$ yang mengatur besarnya langkah pembaruan. Analisis sensitivitas menunjukkan bahwa nilai $a$ yang terlalu besar mengakibatkan osilasi energi yang ekstrem, yang menghalangi sirkuit untuk mencapai titik minimum yang presisi. Di sisi lain, nilai $c$ yang terlalu kecil membuat estimasi gradien menjadi sangat sensitif terhadap \textit{noise} pengukuran, yang menyebabkan algoritma mudah terjebak pada minimum lokal semu. Melalui pengujian pra-simulasi, parameter diatur sedemikian rupa guna memastikan peluruhan (\textit{decaying}) langkah optimasi yang optimal seiring bertambahnya iterasi.

\subsection{Simulasi Satu Langkah Pembaruan Parameter Sudut}
\label{subsec:4_5_3_numerik}

Sebagai ilustrasi numerik, Tabel \ref{tab:vqe_comparison} di Lampiran \ref{appendix:numerical_data} menyajikan riwayat konvergensi energi lintas jendela waktu. Pada iterasi awal, sirkuit mengalami fluktuasi energi yang tinggi sebagai representasi fase eksplorasi parameter. Sebagai contoh, pada jendela waktu tertentu, SPSA berhasil menurunkan nilai energi dari $\langle H \rangle = 0,5$ menjadi $-0,18$ dalam 400 iterasi pertama melalui pembaruan sudut rotasi secara bertahap. Penurunan asimtotik ini membuktikan bahwa \textit{optimizer} berhasil menuntun sistem menuju konfigurasi portofolio dengan imbal hasil maksimal dan risiko minimal sesuai dengan target fisis Hamiltonian.
